\documentclass[11pt]{article}

\usepackage{amsmath, amssymb}
\usepackage{geometry}
\usepackage{natbib}
\usepackage{hyperref}

\geometry{margin=1in}

\title{More Structural Shocks than Observables in DSGE Estimation: A Short Technical Note}
\author{H. Hollander}
\date{}

\begin{document}

\maketitle

\section{State-Space Representation}

A linear(ized) DSGE model can be written in state-space form:
\begin{align}
x_t &= F(\theta)\, x_{t-1} + G(\theta)\, \varepsilon_t, 
& \varepsilon_t \sim \mathcal{N}(0,Q), \label{eq:transition}\\
y_t &= H(\theta)\, x_t + u_t,
& u_t \sim \mathcal{N}(0,R), \label{eq:measurement}
\end{align}
where $x_t$ is a vector of latent states, $y_t$ are observables, and $\varepsilon_t$ are structural shocks.
Shocks affect observables indirectly through the state vector and the model's propagation mechanism; thus counting
arguments based on the number of shocks and observables are not, by themselves, decisive. See \citet{an2007},
\citet{herbst2015}, \citet{durbin2012}.

\section{The Kalman Filter and the Likelihood}

The likelihood integrates out the latent state path:
\begin{align}
p(y_{1:T}\mid \theta)
&= \int p(y_{1:T},x_{1:T}\mid \theta)\, dx_{1:T} \nonumber\\
&= \int \left[\prod_{t=1}^T p(y_t\mid x_t,\theta)\right] p(x_{1:T}\mid \theta)\, dx_{1:T}. \label{eq:likelihood_integral}
\end{align}
Under Gaussian assumptions, the Kalman filter yields the prediction-error decomposition:
\begin{equation}
\log L(\theta)
= \sum_{t=1}^{T} \left[
-\frac{1}{2}\log|S_t|
-\frac{1}{2}\nu_t' S_t^{-1}\nu_t
\right] + \text{constant}, \label{eq:kalman_ll}
\end{equation}
where $\nu_t = y_t - \mathbb{E}[y_t \mid y_{1:t-1}]$ is the one-step-ahead forecast error and
$S_t = \text{Var}(\nu_t \mid y_{1:t-1})$ its variance. This formulation makes clear that the likelihood is well-defined without a
one-to-one mapping from shocks to observables; identification instead relies on the model's dynamic structure and
cross-equation restrictions. See \citet{durbin2012}, \citet{hamilton1994}, \citet{an2007}.

\paragraph{Useful explainers.}
Wikipedia (Kalman filter): \url{https://en.wikipedia.org/wiki/Kalman_filter}.
Standard technical references: \citet{durbin2012} and \citet{hamilton1994}. A DSGE-focused state-space primer is \citet{an2007}.

\section{Why Medium-Scale NK Models Can Sustain Many Shocks}

In medium-scale NK models (e.g., \citealt{smets2007}), different shocks enter different blocks of the model (Euler equations,
Phillips curves, capital accumulation, policy rules) and propagate through common state dynamics. Identification typically
derives from distinct persistence and distinct dynamic responses (IRFs) across observables.

\section{When ``Too Many Shocks'' Becomes a Problem}

Difficulties arise when shocks are observationally equivalent (similar entry points, similar persistence, similar propagation),
leading to weak identification and ``variance swapping.'' A separate concern is that shock flexibility may absorb structural
misspecification and distort inference about structural parameters. See \citet{denhaan2021asd}. Correlated or dynamically
interacting disturbances can also matter empirically; see \citet{curdia2012}.

\section{Identification Criterion}

The relevant condition is not:
\begin{equation}
\#\text{shocks} \le \#\text{observables}.
\end{equation}
Instead, identification depends on rank / local invertibility conditions: whether the mapping from structural parameters (and
shock process parameters) into the likelihood-implied moments (or directly into the state-space likelihood) is locally one-to-one.
See \citet{iskrev2010} for a formal treatment of local identification in DSGE models.

\section{Post-2008 Financial-Friction Models}

Post-2008 DSGE models commonly add financial wedges/shocks (risk premia, credit spreads, intermediation constraints,
capital quality/valuation, liquidity, and sometimes housing/collateral constraints) to match spread dynamics and macro-financial
comovement. This added flexibility can improve empirical fit but increases the risk of observational equivalence among shocks
and of absorbing misspecification. Good practice includes reporting variance decompositions and IRFs, checking identification
diagnostics, and testing shock specifications (cf. \citealt{denhaan2021asd}; \citealt{curdia2012}).

Representative references include \citet{bgg1999}, \citet{iacoviello2005}, \citet{gerali2010}, \citet{gertlerkaradi2011},
\citet{christiano2014}.

\section*{Conclusion}

More shocks than observables is not inherently problematic in DSGE estimation. The likelihood is computed by integrating
over latent states, and identification comes from dynamic structure and cross-equation restrictions. Problems arise when shocks
generate indistinguishable dynamics or absorb misspecification---risks that are especially salient in macro-financial models.

\bibliographystyle{apalike}
\begin{thebibliography}{}

\bibitem[An and Schorfheide(2007)]{an2007}
An, S. and Schorfheide, F. (2007).
Bayesian analysis of DSGE models.
\textit{Journal of Economic Dynamics and Control}.

\bibitem[Bernanke et al.(1999)]{bgg1999}
Bernanke, B., Gertler, M., and Gilchrist, S. (1999).
The financial accelerator in a quantitative business cycle framework.
\textit{Quarterly Journal of Economics}.

\bibitem[Christiano et al.(2014)]{christiano2014}
Christiano, L., Motto, R., and Rostagno, M. (2014).
Risk shocks.
\textit{American Economic Review}.

\bibitem[Cúrdia and Reis(2012)]{curdia2012}
Cúrdia, V. and Reis, R. (2012).
Correlated disturbances and {U.S.} business cycles.
\textit{Journal of Monetary Economics}.

\bibitem[Den Haan and Drechsel(2021)]{denhaan2021asd}
Den Haan, W. and Drechsel, T. (2021).
Agnostic structural disturbances (ASDs): Detecting and reducing misspecification in empirical macroeconomic models.
\textit{Journal of Monetary Economics}.

\bibitem[Durbin and Koopman(2012)]{durbin2012}
Durbin, J. and Koopman, S. (2012).
\textit{Time Series Analysis by State Space Methods}.
Oxford University Press.

\bibitem[Gerali et al.(2010)]{gerali2010}
Gerali, A., Neri, S., Sessa, L., and Signoretti, F. (2010).
Credit and banking in a DSGE model of the euro area.
\textit{Journal of Money, Credit and Banking}.

\bibitem[Gertler and Karadi(2011)]{gertlerkaradi2011}
Gertler, M. and Karadi, P. (2011).
A model of unconventional monetary policy.
\textit{Journal of Monetary Economics}.

\bibitem[Hamilton(1994)]{hamilton1994}
Hamilton, J. (1994).
\textit{Time Series Analysis}.
Princeton University Press.

\bibitem[Herbst and Schorfheide(2015)]{herbst2015}
Herbst, E. and Schorfheide, F. (2015).
\textit{Bayesian Estimation of DSGE Models}.
Princeton University Press.

\bibitem[Iacoviello(2005)]{iacoviello2005}
Iacoviello, M. (2005).
House prices, borrowing constraints, and monetary policy in the business cycle.
\textit{American Economic Review}.

\bibitem[Iskrev(2010)]{iskrev2010}
Iskrev, N. (2010).
Local identification in DSGE models.
\textit{Journal of Monetary Economics}.

\bibitem[Smets and Wouters(2007)]{smets2007}
Smets, F. and Wouters, R. (2007).
Shocks and frictions in U.S. business cycles: A Bayesian DSGE approach.
\textit{American Economic Review}.

\end{thebibliography}

\end{document}